\documentclass[10pt, aspectratio=169]{beamer}
\usetheme{Madrid}
\usecolortheme{dolphin}
\usepackage{ctex}
\usepackage{xeCJK}
\usepackage{amsmath,amssymb}
\usepackage{listings}
\usepackage{tikz}
\usepackage{xcolor}
\title{408考研零基础 --- 数据结构}
\subtitle{1-18 排序入门}
\author{}
\date{}
\setbeamertemplate{page number in head/foot}{}
\begin{document}
\begin{frame}{本节目标}
\begin{enumerate}
\item 理解排序的基本概念（稳定性、内/外部排序）
\item 掌握插入排序、冒泡排序、简单选择排序
\item 理解快速排序的分治思想
\end{enumerate}
\end{frame}
\begin{frame}{排序的基本概念}
\textbf{稳定性}：相等元素排序后相对顺序不变
\vspace{0.3em}
\textbf{分类}
\begin{itemize}
\item 内部排序：所有元素在内存中
\item 外部排序：元素需借助外存
\end{itemize}
\end{frame}
\begin{frame}{插入排序}
\textbf{思想}：将元素插入已排序区的合适位置
\vspace{0.3em}
\textbf{性能}
\begin{itemize}
\item 最好 $O(n)$（已有序）
\item 最坏 $O(n^2)$（逆序）
\item 平均 $O(n^2)$
\end{itemize}
\vspace{0.3em}
空间 $O(1)$，\textbf{稳定}
\end{frame}
\begin{frame}{冒泡排序}
\textbf{思想}：相邻比较交换，最大元素移至末尾
\vspace{0.3em}
\textbf{性能}
\begin{itemize}
\item 最好 $O(n)$（已有序，提前结束）
\item 最坏 $O(n^2)$（逆序）
\item 平均 $O(n^2)$
\end{itemize}
\vspace{0.3em}
空间 $O(1)$，\textbf{稳定}
\end{frame}
\begin{frame}{简单选择排序}
\textbf{思想}：每趟选最小元素与待排序区第一个交换
\vspace{0.3em}
\textbf{性能}
\begin{itemize}
\item 始终 $O(n^2)$（与初始状态无关）
\item 空间 $O(1)$
\end{itemize}
\textbf{不稳定}
\end{frame}
\begin{frame}{快速排序}
\textbf{思想}：分治 + 分区
\vspace{0.3em}
\textbf{分区操作}
\begin{itemize}
\item 选枢轴（通常第一个元素）
\item low/high 指针交替扫描交换
\item 左边 $\le$ 枢轴 $\le$ 右边
\end{itemize}
\vspace{0.3em}
\textbf{性能}
\begin{itemize}
\item 最好 $O(n\log n)$（均匀分割）
\item 最坏 $O(n^2)$（已有序）
\item 平均 $O(n\log n)$
\end{itemize}
空间 $O(\log n)$ 至 $O(n)$，\textbf{不稳定}
\end{frame}
\begin{frame}{排序算法对比}
\begin{tabular}{|c|c|c|c|c|}
\hline 算法 & 最好 & 平均 & 最坏 & 稳定 \\ \hline
插入排序 & $O(n)$ & $O(n^2)$ & $O(n^2)$ & 稳定 \\ \hline
冒泡排序 & $O(n)$ & $O(n^2)$ & $O(n^2)$ & 稳定 \\ \hline
选择排序 & $O(n^2)$ & $O(n^2)$ & $O(n^2)$ & 不稳 \\ \hline
快速排序 & $O(n\log n)$ & $O(n\log n)$ & $O(n^2)$ & 不稳 \\ \hline
\end{tabular}
\vspace{0.5em}
快速排序是实际应用中最常用的排序算法
\end{frame}
\begin{frame}{本节总结}
\begin{enumerate}
\item 插入排序：插入已排序区，稳定
\item 冒泡排序：相邻交换，稳定
\item 选择排序：每趟选最小，不稳定
\item 快速排序：分治 + 分区，最常用
\item 数据结构模块至此结束，共 18 集
\end{enumerate}
\vspace{1em}
\begin{center}
\textcolor{blue}{\textbf{下个模块：计算机组成原理}}
\end{center}
\end{frame}
\end{document}
